001 & Pilot & Direct causal decisions from exact counts. The low-precision arm exhausted its token budget. A format-sensitive endpoint limited the proposed compression comparison. \\[3pt]
002 & Control & Giving the arithmetic formula separated the sign of a numerical estimate from the explicitly selected action. Single-fence handling mattered to the result. \\[3pt]
003 & Technical diagnostic & Disabling CPU repacking did not rescue Q2 output. Repeated items checked operational reproducibility and added no independent families. \\[3pt]
004 & Reference screen & The larger F16 reference missed the competence gate on direct prompts. That gate constrained interpretation of a weaker format. \\[3pt]
005 & Development & Supplying randomized and historical rates made source selection explicit. Action preference and tree-model misfit complicated the conflict-penalty interpretation. \\[3pt]
006 & Control follow-up & Q8 reproduced the strong F16 sign performance on the reused formula tasks. Signs and actions remained separate endpoints. \\[3pt]
007 & Fresh-family replication & The sign result recurred under the supplied formula. On the same 48 items, F16/Q4 sign scores were 46/31 correct and overt action scores were 24/28 correct. All replies were fenced. \\[3pt]
008 & Development crossover & An interrupted F16 run was retained. Completed order arms, rather than partial-run scores with unattempted items, supplied the order comparison. \\[3pt]
009 & Development crossover & Interleaving both evidence orders reduced the between-run confound. An early summary halved success and activation counts. Corrected counts preserve paired effects. \\[3pt]
010 & Fresh-family follow-up & The 3B order pattern recurred on new numerical families. The family-bootstrap analysis followed earlier observations. The original count underreporting was corrected. \\[3pt]
011 & Development & Historical strength was matched across agreement conditions. The quantized model's action preference could reduce an apparent conflict penalty. Corrected count output is used. \\[3pt]
012 & Reference screen & The matched-strength 7B F16 reference failed conflict-specific competence in both evidence orders under the fence-tolerant score. \\[3pt]
013 & Development crossover & The causal-context frame produced capped continuations. A later prefix diagnostic exposed a parseable opening answer but did not rescue the frozen endpoint. \\[3pt]
014 & Development crossover & A stricter response instruction restored valid outputs. Neither reference framing cleared all performance gates. \\[3pt]
015 & Reference screen & The fresh-family F16 gate passed, enabling a conditional Q4 arm. Its observed conflict-order pattern motivated a separate prospective hypothesis. \\[3pt]
016 & Prospective confirmation & A frozen 96-family design confirmed the Q4 conflict-order contrast. Precision interactions were secondary. The magnitude counterexample is retrospective. \\[3pt]
017 & Exploratory extension & Larger 14B quantized arms reused the 016 fixture. The large 7B conflict-order effect did not recur. Archived responses supported post-run analysis recovery. \\[3pt]
018 & Development gate & Matched controls and graded evidence were introduced. F16 failed the decision gate, so no Q4 comparison followed. \\[3pt]
019 & Development gate & Semantic and reversed neutral response contracts changed the mapping to an action. Conflicting prose about which source came first weakened a comprehension probe. The official gate failed. \\[3pt]
020 & Component screen & Single-number and randomized-only components were comparatively tractable. Adding a second source exposed interference and order dependence. \\[3pt]
021 & Guidance screen & All 128 governing-source selections were correct. Ordinary and explicitly guided two-source decisions each scored 58/128. Naming evidence and acting on it are distinct outputs. \\[3pt]
022 & Development replication & A fresh matched fixture compared Qwen and Mistral deployments. Failure on simpler Mistral components limited cross-checkpoint attribution. \\[3pt]
023 & Measurement audit & Crossing source names, physical display and option position left a two-source deficit. Weak harmful-direction numeric and noncausal controls precluded specificity. \\[3pt]
024 & Development matrix & Three lineages and seven formats completed the matrix. Output caps and first-option policies affected particular arms. No F16 passed all of the later component/control diagnostics. \\[3pt]
025 & Retrospective & Archived responses were reanalyzed for matched transitions, information dependence and metamorphic behavior. Reused observations cannot confirm an independently selected hypothesis. \\[3pt]
025.1 & Retrospective audit & An applicability audit reproduced the archive screen and distinguished methods the existing data could support from methods requiring new graded evidence or interventions. \\[3pt]
026 & Development gate & Observationally identical confounder/mediator graphs required different decisions. A faulty complemented-outcome arithmetic key invalidated the first pilot. Repaired forms failed qualification. \\[3pt]
027 & Development gate & Fully crossing answer order addressed a position confound. No tested development arm qualified. Higher-memory models and later validation/precision arms contributed no outcome data. \\[3pt]
028 & Development gate & Bare, structured and calculator-assisted procedures separated computation from procedure choice. Development finished and the gate stopped the pipeline. The earlier pause in the README is stale. \\[3pt]
029 & Prospective case study & CF-TP1 crossed feasible provenance, target mix, order, form and matched arithmetic. The primary contrast was inconclusive. No new precision arm was part of this protocol. \\[3pt]
